\documentclass[12pt]{article}
\usepackage[a4paper, left = 2cm, right = 2cm, tmargin = 0.8cm, bmargin = 0.8cm]{geometry}
\usepackage{xargs}
\usepackage{amsmath,amssymb}
\usepackage{hyperref}
\usepackage{physics}
\usepackage{graphicx}

\usepackage{enumitem}

\setlist[enumerate]{
    topsep=2pt,
    itemsep=1pt,
    parsep=0pt,
    partopsep=0pt
}

\setlist[enumerate,1]{label=\arabic*}
\setlist[enumerate,2]{label=\theenumi.\arabic*}
\setlist[enumerate,3]{label=\theenumii.\arabic*}


\begin{document}

\thispagestyle{empty}

	\begin{center}
	\textbf{Thesis plan:} Coherent states in numerical methods for fermionic particle-preserving quantum systems (Michal Horanský)
	\end{center}
	
	\begin{enumerate}
	
	\item Introduction
	
	\begin{enumerate}
	
		\item Semi-classical methods in high-dimensional quantum mechanics
		
		\begin{enumerate}
		
			\item The problem of diagonalisation
			
			\item The problem of dynamics
		
		\end{enumerate}
		
		\item Coherent states throughout the 20th century
		
		\begin{enumerate}
		
			\item Schr\"odinger coherent states
			
			\item Glauber coherent states
			
			\item Gilmore-Perelomov coherent states
		
		\end{enumerate}
	
	\end{enumerate}
	
	\item Constructing the coherent state manifold
	
	\begin{enumerate}
	
		\item Gilmore-Perelomov coherent states in general
		
		\begin{enumerate}
		
			\item Action group, stability group
	
			\item Reference state operator and Hadamard's lemma
	
			\item Topology of the coherent state manifold
	
			\item Equations of motion
		
		\end{enumerate}
		
		\item $SU(M)$ fermionic coherent states
		
		\begin{enumerate}
		
			\item Displacement operator
	
			\item Decomposition into the full occupancy basis
	
			\item Overlaps and normalisation
	
			\item Fermionic operator sequence matrix elements
	
			\item Overlap update formalism
		
		\end{enumerate}
		
		\item Approximately-coherent particle-number-conserving states
		
		\begin{enumerate}
		
			\item Qubit operators
			
			\item Overlaps with elementary symmetric polynomial
			
			\item Fermionic operator sequence matrix elements
			
			\item Topology of the ACPNCS manifold
		
		\end{enumerate}
	
	\end{enumerate}
	
	
	\item Molecular electronic structure
	
	\begin{enumerate}
	
		\item System description
		
		\begin{enumerate}
		
			\item Non-relativistic electronic Hamiltonian
			
			\item Spin-$z$-labelled subspaces
			
			\item Symmetry-labelled subspaces
			
			\item Potential energy surfaces
		
		\end{enumerate}
		
		\item Ground state estimation with Monte Carlo
		
		\begin{enumerate}
		
			\item Low-excitation-guided sampling
			
			\item Varying the sample basis size
			
			\item Varying the Hilbert space dimension
			
			\item The effects of spin and symmetry
			
			\item Multireferential regimes
		
		\end{enumerate}
		
		\item Potential energy surface extrapolation with coherent-state transfer
		
		\begin{enumerate}
		
			\item Sampling and decomposition strategy
			
			\item Energy eigenstate reconstruction from transferred basis
			
			\item Spin-protected crossings
			
			\item Symmetry-protected crossings
			
			\item Unprotected crossings
			
			\item Multireferential regimes
		
		\end{enumerate}
		
		\item Discussion
	
	\end{enumerate}
	
	
	\item Appendix
	
	\begin{enumerate}
	
		\item[A] Counting constrained minor products
		
		\item[B] Repeated commutator of the $SU(M)$ displacement operator and the $\hat{\phi}$-operator
	
	\end{enumerate}
	
	\end{enumerate}
	

	
\end{document}